% Plantilla de Presentación Beamer Genérica
\documentclass[aspectratio=169, xcolor={dvipsnames}, 8pt, spanish]{beamer}

\usepackage{booktabs}   
\usepackage{multirow}   
\usepackage{amsmath}
\usepackage{amssymb}
\usepackage{caption}    
\usepackage{array}      
\usepackage{tabularx}   
\usepackage{siunitx}    
\usepackage{hyperref}

% --- AJUSTES Y CONFIGURACIÓN ---
% Si el archivo customizacoes.tex existe, se incluye.
% De lo contrario, definimos configuraciones por defecto.
\input{customizacoes}

% Información del Documento Genérico
\title[[Título Corto]]{[Título Completo de la Presentación]}
\subtitle{[Subtítulo de la Presentación / Proyecto]}
\author[[Autor Corto]]{[Nombre Completo del Expositor]}
\institute[[Siglas Institución]]{[Nombre de la Institución / Universidad]}
\date{\today}

% Nombre del Docente / Supervisor
\newcommand{\advisorname}{\textbf{Supervisor / Evaluador:} [Nombre del Evaluador]}

\begin{document}

% Slide 1: Portada
\begin{frame}
    \titlepage
\end{frame}

% Slide 2: Introducción y Contexto
\begin{frame}{Introducción y Justificación}

\begin{columns}[T,onlytextwidth]
\begin{column}{0.32\textwidth}
\begin{block}{1. Contexto}
    Presentación breve del contexto del problema o del área de estudio. Justificar la relevancia del proyecto.
\end{block}
\end{column}
\begin{column}{0.32\textwidth}
\begin{block}{2. Motivación}
    ¿Por qué es necesario este estudio? Identificar la brecha de conocimiento o el problema práctico a resolver.
\end{block}
\end{column}
\begin{column}{0.32\textwidth}
\begin{block}{3. Enfoque}
    Resumen breve de la metodología o el enfoque analítico propuesto para dar respuesta al problema.
\end{block}
\end{column}
\end{columns}

\begin{block}{Objetivo General}
    Definir el objetivo central del proyecto de manera clara y concisa en este bloque destacado.
\end{block}

\end{frame}

% Slide 3: Datos e Instrumentos
\begin{frame}{Metodología y Datos Utilizados}

\begin{columns}[T,onlytextwidth]
\begin{column}{0.48\textwidth}
\begin{block}{Descripción de la Muestra}
    \begin{itemize}
        \item \textbf{Tamaño Muestral ($N$):} [Número de observaciones]
        \item \textbf{Criterios de Inclusión:} [Detallar brevemente]
        \item \textbf{Variables de Control:} [Detallar brevemente]
    \end{itemize}
\end{block}
\end{column}
\begin{column}{0.48\textwidth}
\begin{block}{Diccionario Resumido}
    \begin{table}[]
    \scriptsize
    \begin{tabularx}{\textwidth}{lX}
    \toprule
    \textbf{Variable} & \textbf{Descripción / Rol} \\
    \midrule
    \texttt{Var\_Respuesta} & Respuesta primaria de interés. \\
    \texttt{Var\_Indep\_1} & Primer predictor o covariable. \\
    \texttt{Var\_Indep\_2} & Segundo predictor o factor. \\
    \texttt{Var\_Ajuste} & Variable para control de sesgo. \\
    \bottomrule
    \end{tabularx}
    \caption{Definición de variables clave.}
    \end{table}
\end{block}
\end{column}
\end{columns}

\end{frame}

% Slide 4: Resultados
\begin{frame}{Resultados Principales}

\begin{columns}[T,onlytextwidth]
\begin{column}{0.58\textwidth}
\begin{block}{Análisis y Modelamiento}
    Explicación detallada de los resultados del modelo. 
    \begin{itemize}
        \item Los coeficientes muestran una relación directa significativa ($p < 0.05$).
        \item El poder explicativo del modelo de regresión es adecuado.
    \end{itemize}
\end{block}
\end{column}
\begin{column}{0.38\textwidth}
\begin{block}{Representación Visual}
    \begin{center}
        \color{gray}
        [Espacio reservado para Gráficos de Resultados]
    \end{center}
\end{block}
\end{column}
\end{columns}

\end{frame}

% Slide 5: Conclusiones
\begin{frame}{Conclusiones y Recomendaciones}

\begin{itemize}
    \item \textbf{Resultado Clave 1:} Resumen del hallazgo más importante del estudio.
    \item \textbf{Resultado Clave 2:} Conclusión secundaria relevante.
    \item \textbf{Implicancia:} ¿Cómo impacta esto en la práctica o teoría?
    \item \textbf{Líneas Futuras:} Próximos pasos sugeridos para expandir el análisis.
\end{itemize}

\end{frame}

\end{document}